\section{Datos}

\subsection{Tipos de datos y conversión}

\subsubsection{Tipos de datos}

Los tipos de datos que se pueden encontrar son:
\begin{itemize}
    \item int64: enteros.
    \item float64: números decimales.
    \item object: texto.
    \item datetime64: fechas.
\end{itemize}

Para ver los tipos de datos de una \textit{Serie} o de un \textit{Data Frame} se puede utilizar: \texttt{.dtypes}. 

Si quiere ver los tipos de datos ejecute:
\begin{pycode}{Tipos de datos con dtype}
import pandas as pd

print(df.dtypes)
\end{pycode}

Hay otra forma para ver el tipo de datos de nuestra \textit{Serie} o \textit{Data Frame}, está es \texttt{.info()}.

Además \texttt{.info()} nos devuelve un resumen que incluye:
\begin{enumerate}
    \item Clase de objeto (si es o no un objeto de pandas)
    \item Rango del Índice: Cuántas filas tiene la tabla (ej: RangeIndex: 500 entries, 0 to 499).
    \item Conteo de Columnas: Total de columnas.
    \item Detalle por Columna: Nombre, Non-Null Count (Cuántos datos reales hay, es útil para detectar si faltan datos o hay celdas vacías), Dtype: El tipo de dato (número, texto, fecha).
    \item Uso de Memoria: Cuánta memoria RAM está consumiendo ese DataFrame.
\end{enumerate}

\begin{ejemplo}
    Si tuviera una tabla con 50.000 datos, pero solo tiene 48.000 non-null, esto significa que faltan 2.000 valores.
    Si tuviera un \textit{Data Frame} y usted quisiera realizar un cálculo. Si el \textit{Data Frame} tiene un objeto (texto), no podriamos realizar el cálculo.
\end{ejemplo}

Vease en código:
\begin{pycode}{.info()}
import pandas as pd

Supongamos un DataFrame
df.info()

"""
Salida típica:
<class 'pandas.core.frame.DataFrame'>
RangeIndex: 4 entries, 0 to 3
Data columns (total 3 columns):

#   Column  Non-Null Count  Dtype
0   item    4 non-null      object
1   peso    4 non-null      float64
2   fecha   3 non-null      object  <-- ¡Hay menos datos!
dtypes: float64(1), object(2)
memory usage: 224.0+ bytes
"""
\end{pycode}

\subsubsection{Conversión de datos}

La función \texttt{.astype()} es la herramienta de \textit{Pandas} que permite cambiar el tipo de dato (\textit{dtype}) de una \textit{Serie} o de un \textit{DataFrame} completo.

Es fundamental realizar estas conversiones por:
\begin{enumerate}
\item \textit{Cálculos:} No se puede sumar un número si \textit{Pandas} cree que es una palabra (\textit{string}).
\item \textit{Eficiencia:} Algunos tipos de datos consumen mucha menos memoria que otros.
\end{enumerate}

Lost tipos de datos comunes son:
\begin{itemize}
\item \textit{int} / \textit{float}: Para cálculos numéricos.
\item \textit{str} / \textit{object}: Para texto o datos alfanuméricos.
\item \textit{category}: Se utiliza cuando una columna tiene pocos valores que se repiten mucho (ej: nombres de países o tipos de moneda). Al usar \texttt{'category'}, el archivo se vuelve mucho más liviano y rápido de procesar.
\end{itemize}

Veamos cómo aplicarlo de forma individual o masiva:

\begin{pycode}{Uso de .astype()}
import pandas as pd

# Data frame
df = pd.DataFrame({
'ID': ['101', '102', '103'], # Se carga como texto (object)
'Precio': [450.5, 300.0, 150.2],     # Se carga como decimal (float)
'Provincia': ['Buenos Aires', 'Mendoza', 'Cordoba'] # Se carga como texto (object)
})

# Conversión de una columna
df['ID'] = df['ID'].astype(int) # Cambiamos de texto a int

# Conversión múltiple
df = df.astype({
'Precio': int,
'Provincia': 'category'
})

print(df.dtypes)
"""
ID           int64
Precio       int64
Provincia    category
dtype: object
"""
\end{pycode}

\begin{nota}
    Al convertir de \textit{float} (decimal) a \textit{int} (entero), \textit{Pandas} no redondea, borra los decimales.
\end{nota}





\subsection{Manejo de índices y etiquetas}

Por defecto, el índice en \textit{Pandas} es un número entero secuencial (0,1,2,…). Sin embargo, para que nuestros datos tengan sentido, muchas veces necesitamos que las etiquetas de las filas y columnas sean descriptivas.

\subsubsection{Ver índices actuales (\textit{.index} y \textit{.columns})}

La propiedad \texttt{.index} permite inspeccionar cómo se llaman sus filas, mientras que \texttt{.columns} hace lo mismo con los encabezados. 

\begin{pycode}{Inspección de etiquetas}
import pandas as pd

print(df.index) #  Muestra el rango de las filas (ej: de 0 a 100)
print(df.columns) # Muestra una lista con los nombres de las columnas
\end{pycode}

\subsubsection{Renombrar etiquetas (\textit{.rename})}

Para cambiar los nombres de columnas o filas de forma selectiva sin afectar al resto del \textit{DataFrame}, utilizamos la función \texttt{.rename()}. Esta función utiliza diccionarios de \textit{Python} para mapear el nombre viejo con el nuevo.

\begin{verbatim}
df.rename(columns={'Nombre\_Viejo': 'Nombre\_Nuevo'})
\end{verbatim}

\begin{pycode}{Uso de .rename()}
import pandas as pd

df = df.rename(columns={
'Precio': 'Valor_Unitario_USD',
'Qty': 'Cantidad_Bultos'
})

# También se puede aplicar a las filas (índices)
df = df.rename(index={0: 'Enero', 1: 'Febrero'})
\end{pycode}

\subsubsection{Reiniciar el contador del índice (\textit{.reset\_index})}

Cuando usted borra filas o filtra datos, el índice original se desordena (podría quedar una secuencia como 0,5,12,…). Para volver a tener una numeración limpia y consecutiva, usamos \texttt{.reset\_index()}.

\begin{nota}
El parámetro \textit{drop}: Por defecto \texttt{.reset\_index()} convierte el índice viejo en una columna nueva. Si usted no quiere conservar esos números viejos y prefiere borrarlos para siempre, debe usar \texttt{drop=True}.
\end{nota}

\begin{pycode}{Limpieza de índices}
import pandas as pd

# Después de un filtrado, reiniciamos el contador desde 0
df_limpio = df.reset_index(drop=True)
\end{pycode}



\subsection{Selección y filtrado}

Además de la selección posicional o por etiquetas (\textit{.loc} e \textit{.iloc}), \textit{Pandas} tiene otras formas de filtrar datos.

\subsubsection{Filtrado con máscaras booleanas}

Una máscara booleana es una serie de valores \textit{True} o \textit{False} que \textit{Pandas} genera al evaluar una condición en cada fila. Al aplicar esta máscara al \textit{DataFrame}, la librería solo nos muestra las filas donde el resultado fue \textit{True}.

Podemos utilizar operadores de comparación estándar: \texttt{>}, \texttt{<}, \texttt{==} (igual), \texttt{!=} (distinto), \texttt{>=} y \texttt{<=}.

\begin{pycode}{Máscara Booleana}
import pandas as pd

# Creamos la máscara (una serie de True/False)
mask = df["Edad"] > 30

# Aplicamos la máscara al DataFrame para "esconder" los False
df_filtrado = df[mask]
print(df_filtrado)
\end{pycode}

\subsubsection{Filtrado con condiciones lógicas (\textit{AND} y \textit{OR})}

A menudo, una sola condición no es suficiente. Para combinar filtros, utilizamos los operadores lógicos vectorizados de \textit{NumPy}:
\begin{itemize}
\item \textit{\& (AND):} Se deben cumplir ambas condiciones.
\item \textit{| (OR):} Se debe cumplir al menos una de las condiciones.
\end{itemize}

\begin{nota}
\textit{Regla de los paréntesis:} Cuando combine condiciones en \textit{Pandas}, es obligatorio envolver cada condición entre paréntesis \texttt{( )}. Si los olvida, \textit{Python} se confundirá con el orden de las operaciones y le lanzará un error.
\end{nota}

Véalo en código con un ejemplo de salarios y ciudades:

\begin{pycode}{Condiciones lógicas combinadas}
import pandas as pd

# Caso 1: Filtrar empleados Senior con buen sueldo (AND)
senior_bien_pagos = df[(df["Edad"] > 30) & (df["Salario"] > 60000)]

# Caso 2: Filtrar por múltiples ciudades (OR)
norte_y_litoral = df[(df["Ciudad"] == "Rosario") | (df["Ciudad"] == "Salta")]

print(senior_bien_pagos)
print(norte_y_litoral)
\end{pycode}

